%%
%% part3-customization/ch17-options.tex
%%
%% Chapter 17: Class Options — draft/final modes, font
%% selection, and one/two-sided layout.
%%

\chapter{آپشن‌های کلاس}
\label{chap:options}

\section{چرا این موضوع مهم است؟}
\label{sec:options-why}

آپشن‌های کلاس، به شما امکان می‌دهند که
\lr{MatinBook} را برای نیازهای خاص خود تنظیم
کنید. از انتخاب فونت تا حالت پیش‌نویس، همه
با یک خط \cmd{documentclass} قابل کنترل
هستند.

\lr{MatinBook} از سیستم \textbf{کلید-مقدار}
\lr{LaTeX 2023} استفاده می‌کند که بسیار
مدرن‌تر و پایدارتر از روش‌های قدیمی است.
در این فصل، با تمام آپشن‌های کلاس آشنا
می‌شوید.

\mnote{سیستم \lr{\textbackslash DeclareKeys} در \lr{LaTeX 2023} جایگزین \lr{\textbackslash DeclareOption} قدیمی شده و از سینتکس \lr{key=value} پشتیبانی می‌کند.}

\section{آپشن‌های حالت}
\label{sec:options-mode}

\subsection{حالت \lr{final}}
\label{sec:options-final}

حالت پیش‌فرض \lr{MatinBook}، \textbf{حالت
نهایی} است. در این حالت، کتاب به‌صورت
حرفه‌ای و بدون هیچ نشانه‌ی اضافی چاپ
می‌شود.

\begin{latin}
\begin{minted}{latex}
\documentclass[final]{matinbook}
\end{minted}
\end{latin}

\mnote{حالت \lr{final} نیازی به ذکر ندارد، چون پیش‌فرض است. یعنی \lr{\textbackslash documentclass\{matinbook\}} معادل \lr{\textbackslash documentclass[final]\{matinbook\}} است.}

\subsection{حالت \lr{draft}}
\label{sec:options-draft}

در حالت \textbf{پیش‌نویس}، \lr{MatinBook}
رفتارهای زیر را تغییر می‌دهد:

\begin{itemize}
    \item \textbf{واترمارک:} عبارت \lr{DRAFT}
    با شفافیت \pnum{۹۰٪} در پس‌زمینه‌ی هر
    صفحه نمایش داده می‌شود.
    \item \textbf{جعبه‌های سرریز:} خطوطی که از
    حاشیه خارج می‌شوند، با خط مشکی \pnum{۵pt}
    مشخص می‌شوند.
    \item \textbf{تصاویر:} به‌جای نمایش تصاویر،
    فقط کادر آن‌ها رسم می‌شود (سرعت کامپایل
    بالاتر).
\end{itemize}

\begin{latin}
\begin{minted}{latex}
\documentclass[draft]{matinbook}
\end{minted}
\end{latin}

\mnote{حالت \lr{draft} برای مراحل اولیه‌ی نوشتن بسیار مفید است، ولی قبل از چاپ نهایی حتماً به \lr{final} برگردید.}

\section{آپشن‌های فونت}
\label{sec:options-fonts}

\lr{MatinBook} از پنج فونت فارسی پشتیبانی
می‌کند. برای انتخاب فونت، نام آن را به‌عنوان
آپشن به کلاس بدهید:

\begin{latin}
\begin{minted}{latex}
% |\pc{فونت پیش‌فرض}|%
\documentclass[niloofar]{matinbook}

% |\pc{فونت وزیرمتن}|%
\documentclass[vazirmatn]{matinbook}

% |\pc{فونت ساحل}|%
\documentclass[sahel]{matinbook}

% |\pc{فونت ایران‌لاتوس}|%
\documentclass[irlotus]{matinbook}

% |\pc{فونت بی‌نازنین}|%
\documentclass[bnazanin]{matinbook}
\end{minted}
\end{latin}

جدول \cref{tab:font-options} این پنج فونت
را مقایسه می‌کند.

\begin{matintable}{فونت‌های پشتیبانی‌شده‌ی \lr{MatinBook}}{tab:font-options}
\begin{tblr}{
    width = \textwidth,
    colspec = {l X[l] X[c]},
    hlines, vlines,
    colsep = 8pt,
    rowsep = 4pt,
    row{1} = {font=\bfseries},
}
آپشن & فونت & پیشنهاد \\
\lr{niloofar} & \lr{XB Niloofar} & پیش‌فرض \\
\lr{vazirmatn} & \lr{Vazirmatn} & مدرن \\
\lr{sahel} & \lr{Sahel} & ساده \\
\lr{irlotus} & \lr{IR Lotus} & سنتی \\
\lr{bnazanin} & \lr{B Nazanin} & تجاری \\
\end{tblr}
\end{matintable}

\mnote{فونت \lr{B Nazanin} تجاری است و ممکن است روی همه‌ی سیستم‌ها نصب نباشد. برای کتاب‌های فنی، \lr{XB Niloofar} یا \lr{Vazirmatn} توصیه می‌شود.}

\subsection{ترکیب چند آپشن}
\label{sec:options-combine}

می‌توانید چند آپشن را با کاما از هم جدا کنید:

\begin{latin}
\begin{minted}{latex}
\documentclass[draft, vazirmatn]{matinbook}
\end{minted}
\end{latin}

\mnote{ترتیب آپشن‌ها مهم نیست. \lr{[draft, vazirmatn]} معادل \lr{[vazirmatn, draft]} است.}

\section{آپشن‌های چیدمان}
\label{sec:options-layout}

\subsection{حالت دوطرفه و یک‌طرفه}
\label{sec:options-twoside}

\lr{MatinBook} به‌طور پیش‌فرض از حالت
\textbf{دوطرفه} (\lr{twoside}) استفاده
می‌کند. یعنی صفحات فرد و زوج، چیدمان
آینه‌ای دارند (حاشیه‌ی پهن در بیرون).

\begin{latin}
\begin{minted}{latex}
% |\pc{حالت دوطرفه (پیش‌فرض)}|%
\documentclass[twoside]{matinbook}

% |\pc{حالت یک‌طرفه}|%
\documentclass[oneside]{matinbook}
\end{minted}
\end{latin}

\mnote{حالت \lr{twoside} برای کتاب‌های چاپی مناسب است. حالت \lr{oneside} برای اسناد دیجیتال یا فایل‌های خواندنی روی صفحه‌نمایش.}

\subsection{آپشن‌های \lr{geometry}}
\label{sec:options-geometry}

برای تغییر هندسه‌ی صفحه، می‌توانید از
دستور \cmd{geometry} در پیش از
\cmd{begin\{document\}} استفاده کنید:

\begin{latin}
\begin{minted}{latex}
\usepackage{geometry}
\geometry{
    paperwidth=17.8cm,
    paperheight=25.4cm,
    outer=5.9cm,
    inner=1.5cm
}
\end{minted}
\end{latin}

\mnote{برای جزئیات بیشتر درباره‌ی چیدمان، به \cref{chap:layout} مراجعه کنید.}

\section{آپشن‌های سفارشی}
\label{sec:options-custom}

\lr{MatinBook} از آپشن‌های سفارشی نیز پشتیبانی
می‌کند. برای افزودن آپشن جدید، به فایل
\file{matinbook.cls} بروید و از
\cmd{DeclareKeys} استفاده کنید.

\subsection{مثال: افزودن آپشن \lr{print}}
\label{sec:options-custom-example}

فرض کنید می‌خواهید یک آپشن \lr{print} اضافه
کنید که رنگ‌ها را برای چاپ سیاه‌وسفید تنظیم
کند:

\begin{latin}
\begin{minted}{latex}
% |\pc{در فایل}|\lr{matinbook.cls}:
\newif\if@matin@print
\@matin@printfalse

\DeclareKeys{
    print .code = \@matin@printtrue,
}

\ProcessKeyOptions
\end{minted}
\end{latin}

سپس در ماژول‌های دیگر، از این فلگ استفاده
کنید:

\begin{latin}
\begin{minted}{latex}
\if@matin@print
    % |\pc{تنظیمات چاپ سیاه‌وسفید}|%
\fi
\end{minted}
\end{latin}

\mnote{این یک مثال پیشرفته است. برای کارهای معمول، نیازی به افزودن آپشن جدید نیست.}

\section{آپشن‌ها در فایل‌های مختلف}
\label{sec:options-files}

اگر پروژه‌ی شما چند فایل دارد، آپشن‌ها
فقط در \file{main.tex} تعریف می‌شوند:

\begin{latin}
\begin{minted}{latex}
% |\pc{در فایل}|\lr{main.tex}:
\documentclass[draft, vazirmatn]{matinbook}

\begin{document}
\input{chapters/chapter-01}
\end{document}
\end{minted}
\end{latin}

\mnote{فایل‌های فصل (\lr{chapter-01.tex} و...) نباید \lr{\textbackslash documentclass} داشته باشند. آن‌ها فقط محتوا هستند.}

\section{جدول خلاصه‌ی آپشن‌ها}
\label{sec:options-summary-table}

جدول \cref{tab:all-options} تمام آپشن‌های
\lr{MatinBook} را خلاصه می‌کند.

\begin{matintable}{خلاصه‌ی آپشن‌های \lr{MatinBook}}{tab:all-options}
\begin{tblr}{
    width = \textwidth,
    colspec = {l X[l] X[c]},
    hlines, vlines,
    colsep = 8pt,
    rowsep = 4pt,
    row{1} = {font=\bfseries},
}
آپشن & توضیح & پیش‌فرض \\
\lr{draft} & حالت پیش‌نویس & خیر \\
\lr{final} & حالت نهایی & بله \\
\lr{niloofar} & فونت \lr{XB Niloofar} & بله \\
\lr{vazirmatn} & فونت \lr{Vazirmatn} & خیر \\
\lr{sahel} & فونت \lr{Sahel} & خیر \\
\lr{irlotus} & فونت \lr{IR Lotus} & خیر \\
\lr{bnazanin} & فونت \lr{B Nazanin} & خیر \\
\lr{twoside} & چیدمان دوطرفه & بله \\
\lr{oneside} & چیدمان یک‌طرفه & خیر \\
\end{tblr}
\end{matintable}

\section{خطاهای رایج}
\label{sec:options-mistakes}

\subsection{آپشن ناشناخته}
\label{sec:options-mistake-unknown}

\textbf{مشکل:} اگر آپشن ناشناخته‌ای بدهید،
\lr{LaTeX} آن را به کلاس \lr{book} می‌فرستد
و ممکن است خطا بدهد.

\textbf{راه‌حل:} فقط از آپشن‌های مستند استفاده
کنید.

\subsection{ترکیب آپشن‌های متناقض}
\label{sec:options-mistake-conflict}

\textbf{مشکل:} اگر \lr{draft} و \lr{final} را
با هم بدهید، آخرین آپشن اعمال می‌شود.

\textbf{راه‌حل:} فقط یکی از آپشن‌های متناقض
را بدهید.

\subsection{تعریف آپشن در فایل فصل}
\label{sec:options-mistake-chapter}

\textbf{مشکل:} اگر \cmd{documentclass} را در
فایل فصل بنویسید، خطا می‌دهد.

\textbf{راه‌حل:} \cmd{documentclass} فقط در
\file{main.tex}.

\section{تمرین‌ها}
\label{sec:options-exercises}

\begin{exercise}
    کتاب خود را در حالت \lr{draft} کامپایل
    کنید و واترمارک را ببینید.
    \label{exc:options-draft}
\end{exercise}

\begin{exercise}
    فونت را به \lr{vazirmatn} تغییر دهید و
    نتیجه را مقایسه کنید.
    \label{exc:options-font}
\end{exercise}

\begin{exercise}
    یک کتاب با حالت \lr{oneside} بسازید و
    تفاوت آن را با \lr{twoside} ببینید.
    \label{exc:options-side}
\end{exercise}

\begin{exercise}
    یک آپشن سفارشی (مثلاً \lr{print}) به
    \file{matinbook.cls} اضافه کنید.
    \label{exc:options-custom}
\end{exercise}

\section{خلاصه}
\label{sec:options-summary}

در این فصل، با آپشن‌های کلاس \lr{MatinBook}
آشنا شدیم:

\begin{itemize}
    \item \lr{MatinBook} از سیستم کلید-مقدار
    \lr{LaTeX 2023} استفاده می‌کند.

    \item حالت \lr{final} پیش‌فرض است و
    کتاب را بدون نشانه‌ی اضافی چاپ می‌کند.

    \item حالت \lr{draft} واترمارک، جعبه‌های
    سرریز و تصاویر ساده‌شده را فعال می‌کند.

    \item پنج فونت فارسی پشتیبانی می‌شود:
    \lr{niloofar}، \lr{vazirmatn}، \lr{sahel}،
    \lr{irlotus}، \lr{bnazanin}.

    \item حالت \lr{twoside} پیش‌فرض است و
    برای کتاب‌های چاپی مناسب است.

    \item آپشن‌های سفارشی با
    \cmd{DeclareKeys} قابل افزودن هستند.

    \item آپشن‌ها فقط در \file{main.tex}
    تعریف می‌شوند.

    \item خطاهای رایج شامل آپشن ناشناخته،
    ترکیب آپشن‌های متناقض، و تعریف آپشن در
    فایل فصل است.
\end{itemize}

در فصل بعدی (\cref{chap:colors})، با سیستم
رنگ \lr{MatinBook} آشنا می‌شوید و یاد
می‌گیرید که چگونه رنگ‌ها را شخصی‌سازی
کنید.